\begin{abstract}
We adapt ViT-B/16 to a fine-grained dataset with 10,000 classes and only 45 training images per class, with 2 hours on 48\,GB of GPU memory per training run.
We compare seven adaptation methods under the same settings, from a linear probe to full fine-tuning, BitFit, Adapter, AdaptFormer and LoRA.
The choice that mattered most was how to start the new 10,000-class head.
The usual start, a linear probe trained on frozen features of the training images, made all three methods we tested worse; starting each class from its mean feature (a prototype) gave 2.1--3.1 points more.
We also find that where the trainable parameters are matters more than how many there are: LoRA on 1.375\% of the backbone matches fine-tuning the last four blocks (33\%).
Our final full fine-tuning model reaches 76.22\% top-1 on the CodaBench validation set and 76.09\% on the test set (reference: 70.8\%) in 1.03 GPU-hours, and our final LoRA model reaches 71.90\% and 71.67\% (reference LoRA: 68.3\%) with no extra inference cost.
We also report the ideas that failed, and show after the deadline that larger input images add another 4.4 points.
\end{abstract}
